\chapter{Data Structures --- Trees, Binary Trees, and B+ Trees}
\label{ch:trees_binary_and_bplus}

This chapter covers tree terminology, fundamental tree properties, binary trees, binary search trees, balanced trees, heaps, B-trees, B+ trees, and the mathematical calculations commonly encountered in technical recruitment examinations.

Special attention is given to \textbf{calculating the maximum and minimum number of keys in a B+ tree from its order and height}, as the formulation strictly depends on the architectural definitions employed in the question.

\begin{center}
\vspace{0.5em}
{\color{primary}\rule{0.45\textwidth}{1.2pt}}\\[4pt]
{\Large\bfseries\color{primary} Section I: Fundamentals of Trees}\\[2pt]
{\color{primary}\rule{0.45\textwidth}{1.2pt}}
\vspace{0.5em}
\end{center}

\section{What is a Tree?}

A \textbf{tree} is a hierarchical, non-linear data structure consisting of nodes connected by directed or undirected edges. In a rooted tree, exactly one node is designated as the \textbf{root}, and every other node has exactly one parent.

From a graph-theoretic perspective, a tree is a \textbf{connected, acyclic undirected graph}. This implies that between any two nodes there exists exactly one simple path, and no closed cycles exist.

Trees are indispensable when modeling hierarchical or branching relationships, such as file system folder structures, organizational hierarchies, database indexing (B/B+ trees), Document Object Models (DOM), and compiler abstract syntax trees (AST).

\subsection{Basic Terminology}

\begin{table}[htbp]
\centering
\small
\begin{tabularx}{\textwidth}{l p{7.5cm} X}
\toprule
\textbf{Term} & \textbf{Formal Definition} & \textbf{Example} \\
\midrule
Node             & An individual data element in the tree                      & A directory or file \\
Root             & The topmost node, possessing no parent                      & The root directory (\texttt{/}) \\
Edge             & A directed/undirected link connecting two nodes             & Parent-child relationship \\
Parent           & The immediate predecessor of a node                         & Directory containing a subdirectory \\
Child            & A node directly succeeding another node                     & Subdirectory inside a folder \\
Siblings         & Nodes sharing the identical parent                          & Files in the same directory \\
Leaf Node        & An external node having zero children                       & An empty file / leaf directory \\
Internal Node    & A non-leaf node possessing at least one child               & A folder containing items \\
Ancestor         & Any node along the upward path from a node to the root      & Parent, grandparent, etc. \\
Descendant       & Any node situated in the subtree below a given node         & Child, grandchild, etc. \\
Degree of a Node & Total number of children belonging to the node              & Node with 3 children has degree 3 \\
Degree of a Tree & The maximum degree of any node within the tree              & A binary tree has degree $\leq 2$ \\
Path             & A continuous sequence of connected nodes                    & Root-to-leaf traversal \\
Path Length      & Number of edges traversed along the path                    & 4 edges = path length 4 \\
Depth of a Node  & Number of edges from the root to that node                  & Root depth is $0$ \\
Height of a Node & Number of edges on the longest downward path to a leaf      & A leaf node has height $0$ \\
Height of a Tree & The height of its root node (maximum depth among all nodes) & Edge count on longest path \\
\bottomrule
\end{tabularx}
\caption{Fundamental tree terminology and definitions.}
\end{table}

\begin{quote}
\textbf{Standard Height Convention:} Throughout this text, \textbf{height is measured in edges}. Consequently, a single-node tree has height $0$, and a tree of height $h$ contains $h + 1$ levels (levels $0, 1, \dots, h$).
\end{quote}

\subsection{Fundamental Mathematical Properties of Trees}

\subsubsection{Property 1: Number of Edges}
A tree containing $n$ nodes has exactly:
\[
\boxed{E = n - 1}
\]
where $E$ is the total edge count.
\begin{itemize}
    \item $1$ node $\implies 0$ edges.
    \item $10$ nodes $\implies 9$ edges.
    \item $100$ nodes $\implies 99$ edges.
\end{itemize}
\textbf{Proof:} The root node has zero incoming edges. Each of the remaining $n - 1$ non-root nodes possesses exactly one incoming edge from its unique parent. Hence, total edges $= n - 1$.

\subsubsection{Property 2: Unique Simple Path}
Between any two arbitrary nodes in a tree, there exists \textbf{exactly one simple path}. If two distinct paths connected the same pair, their union would form a cycle, violating the acyclic definition.

\subsubsection{Property 3: Invariance of Edge Count to Topology}
A tree with $n$ nodes contains $n - 1$ edges \emph{regardless} of whether it is perfectly balanced or degenerate (skewed into a linked list). Structural topology governs tree height and query complexity, but never the edge count.

\subsubsection{Property 4: Edge Addition}
Adding an edge between any two existing nodes in a tree creates exactly one cycle.

\subsubsection{Property 5: Edge Removal}
Removing any single edge from a tree partitions it into exactly two disconnected components (subtrees).

\section{General Tree Mathematics}

For general trees where internal nodes can have up to $m$ children ($m$-ary trees):

\subsection{Maximum Nodes at Level $i$}
In a rooted tree of maximum degree $m$, the maximum node capacity at level $i$ (with root at level $0$) is:
\[
\boxed{N_i \leq m^i}
\]
Level $0$ has $m^0 = 1$ node; Level $1$ has $m^1 = m$; Level $2$ has $m^2$; Level $i$ has $m^i$.

\subsection{Maximum Total Nodes for Height $h$}
Summing the geometric progression from level $0$ to level $h$:
\[
N_{\max} = 1 + m + m^2 + \dots + m^h = \boxed{\frac{m^{h+1} - 1}{m - 1}} \quad (m > 1)
\]
For example, for a ternary tree ($m = 3$) of height $h = 3$:
\[
N_{\max} = \frac{3^{3+1} - 1}{3 - 1} = \frac{81 - 1}{2} = 40 \text{ nodes}
\]

\subsection{Minimum Possible Height for $n$ Nodes}
To minimize height, all levels must be packed maximally. Solving $n \leq \frac{m^{h+1}-1}{m-1}$ for $h$:
\[
\boxed{h_{\min} = \left\lceil \log_m\big(n(m - 1) + 1\big) \right\rceil - 1}
\]
For a binary tree ($m = 2$) with $n = 100$:
\[
h_{\min} = \lceil \log_2(101) \rceil - 1 = 7 - 1 = 6 \text{ edges}
\]

\subsection{Maximum Possible Height for $n$ Nodes}
Without structural balancing, a tree can degenerate into a linear chain:
\[
\boxed{h_{\max} = n - 1}
\]

\section{Traversing a Tree}

Tree traversal systematically visits every node in a defined sequence. The four standard traversal methods are:

\subsection{Preorder Traversal (Root $\to$ Left $\to$ Right)}
\begin{lstlisting}[style=pseudocode]
PREORDER(node):
    if node is not null:
        visit(node)
        PREORDER(node.left)
        PREORDER(node.right)
\end{lstlisting}
Used for tree cloning, serialization, and prefix expression evaluation.

\subsection{Inorder Traversal (Left $\to$ Root $\to$ Right)}
\begin{lstlisting}[style=pseudocode]
INORDER(node):
    if node is not null:
        INORDER(node.left)
        visit(node)
        INORDER(node.right)
\end{lstlisting}
\textbf{Key Rule:} In a Binary Search Tree (BST), an inorder traversal visits keys in \textbf{strictly ascending sorted order}.

\subsection{Postorder Traversal (Left $\to$ Right $\to$ Root)}
\begin{lstlisting}[style=pseudocode]
POSTORDER(node):
    if node is not null:
        POSTORDER(node.left)
        POSTORDER(node.right)
        visit(node)
\end{lstlisting}
Used for bottom-up node deallocation (deleting a tree) and postfix (Reverse Polish) expression evaluation.

\subsection{Level-Order Traversal (Breadth-First)}
Visits nodes level-by-level from left to right using a FIFO queue. Used in BFS, tree serialization, and finding shortest root-to-leaf paths.

\textbf{Traversal Complexity:} All four traversals visit every node exactly once:
\[
\text{Time Complexity} = O(n)
\]
Auxiliary space is $O(h)$ for recursive call stacks ($O(n)$ worst-case) and $O(n)$ for level-order queue storage.

\begin{center}
\vspace{1em}
{\color{primary}\rule{0.45\textwidth}{1.2pt}}\\[4pt]
{\Large\bfseries\color{primary} Section II: Binary Trees}\\[2pt]
{\color{primary}\rule{0.45\textwidth}{1.2pt}}
\vspace{0.5em}
\end{center}

\section{What Is a Binary Tree?}

A \textbf{binary tree} is a rooted tree where every node has \textbf{at most two children}, distinguished as the \textbf{left child} and \textbf{right child}. A left child is distinct from a right child even if a parent has only one child.

\subsection{Types of Binary Trees}

\begin{enumerate}
    \item \textbf{Full Binary Tree (Strict / Proper):} Every node has either $0$ or exactly $2$ children. No node has $1$ child.
    \[
    L = I + 1, \quad n = 2I + 1, \quad I = \frac{n - 1}{2}, \quad L = \frac{n + 1}{2}
    \]
    A full binary tree \textbf{must always possess an odd number of nodes}.
    \item \textbf{Complete Binary Tree:} All levels are completely full except possibly the last level, which is filled strictly from left to right without gaps. Used in binary heaps.
    \item \textbf{Perfect Binary Tree:} All internal nodes have $2$ children and all leaves reside at the identical depth. For height $h$:
    \[
    n = 2^{h+1} - 1, \quad L = 2^h, \quad I = 2^h - 1
    \]
    \item \textbf{Balanced Binary Tree:} A binary tree whose height is bounded to $O(\log n)$ to guarantee efficient operations (e.g., AVL, Red-Black).
    \item \textbf{Skewed (Degenerate) Binary Tree:} Every internal node has exactly $1$ child, forming a linear list of height $h = n - 1$.
\end{enumerate}

\subsection{Full vs Complete vs Perfect Comparison}

\begin{table}[htbp]
\centering
\small
\begin{tabularx}{\textwidth}{l X X X}
\toprule
\textbf{Property} & \textbf{Full Binary Tree} & \textbf{Complete Binary Tree} & \textbf{Perfect Binary Tree} \\
\midrule
Internal node degree    & Exactly 2 children        & May have 1 child (at most one) & Exactly 2 children \\
All levels filled       & Not necessarily           & All except possibly last      & Completely filled \\
Last level layout       & No restriction            & Left-to-right without gaps    & Full left-to-right \\
All leaves equal depth  & Not necessarily           & Not necessarily               & Yes \\
Total node parity       & Always odd ($2I+1$)       & Odd or Even                   & Always odd ($2^{h+1}-1$) \\
Node count at height $h$& Varies                    & $2^h \dots 2^{h+1}-1$         & Exactly $2^{h+1}-1$ \\
\bottomrule
\end{tabularx}
\caption{Detailed comparison between Full, Complete, and Perfect binary trees.}
\end{table}

\section{Important Binary Tree Formulas}

Assuming height $h$ is measured in edges and root is at level $0$:

\subsection{Fundamental Extremum Bounds}
\begin{itemize}
    \item \textbf{Max nodes at level $i$:} $N_i^{\max} = 2^i$.
    \item \textbf{Max total nodes for height $h$:} $N_{\max} = 2^{h+1} - 1$.
    \item \textbf{Max leaves for height $h$:} $L_{\max} = 2^h$.
    \item \textbf{Min nodes for height $h$:} $N_{\min} = h + 1$ (skewed tree).
    \item \textbf{Min height for $n$ nodes:} $h_{\min} = \lceil \log_2(n + 1) \rceil - 1$.
\end{itemize}

\subsection{Theorem: Leaves vs Two-Child Nodes ($n_0 = n_2 + 1$)}
In \textbf{any} nonempty binary tree (not just full binary trees), let $n_0$ be leaves, $n_1$ be single-child nodes, and $n_2$ be two-child nodes:
\[
\boxed{n_0 = n_2 + 1}
\]
\textbf{Derivation:}
The total number of nodes is $n = n_0 + n_1 + n_2$.
The total number of edges is $E = n - 1 = n_1 + 2n_2$.
Substituting $n$:
\[
(n_0 + n_1 + n_2) - 1 = n_1 + 2n_2 \implies n_0 - 1 = n_2 \implies \boxed{n_0 = n_2 + 1}
\]

\subsection{Number of Null Child Pointers}
A binary tree with $n$ nodes has $2n$ total pointer fields. Exactly $n - 1$ fields point to child nodes (edges). Hence:
\[
\boxed{\text{Null Pointers} = 2n - (n - 1) = n + 1}
\]

\section{Array Representation of a Binary Tree}

\subsection{Zero-Based Indexing Formulas}
Root at index $0$:
\[
\text{Left}(i) = 2i + 1, \quad \text{Right}(i) = 2i + 2, \quad \text{Parent}(i) = \left\lfloor \frac{i - 1}{2} \right\rfloor
\]

\subsection{One-Based Indexing Formulas}
Root at position $1$:
\[
\text{Left}(i) = 2i, \quad \text{Right}(i) = 2i + 1, \quad \text{Parent}(i) = \left\lfloor \frac{i}{2} \right\rfloor
\]

\section{Number of Structurally Different Binary Trees}

For $n$ unlabeled nodes, the number of distinct binary tree shapes is given by the \textbf{Catalan Number} $C_n$:
\[
\boxed{C_n = \frac{1}{n+1} \binom{2n}{n} = \frac{(2n)!}{(n+1)! \, n!}}
\]
Values: $C_0=1, \; C_1=1, \; C_2=2, \; C_3=5, \; C_4=14, \; C_5=42, \; C_6=132$.
This also equals the number of distinct Binary Search Trees (BSTs) that can be formed from $n$ distinct keys.

\begin{center}
\vspace{1em}
{\color{primary}\rule{0.45\textwidth}{1.2pt}}\\[4pt]
{\Large\bfseries\color{primary} Section III: Binary Search Trees \& Balanced Trees}\\[2pt]
{\color{primary}\rule{0.45\textwidth}{1.2pt}}
\vspace{0.5em}
\end{center}

\section{Binary Search Tree (BST)}

A \textbf{Binary Search Tree} enforces the search ordering property:
\[
\text{All keys in Left Subtree} < \text{Parent Key} < \text{All keys in Right Subtree}
\]

\subsection{Time Complexity}
Search, insertion, and deletion operate in $O(h)$ time:
\begin{itemize}
    \item Balanced BST: $h = O(\log n) \implies O(\log n)$ operations.
    \item Skewed BST: $h = O(n) \implies O(n)$ worst-case operations.
\end{itemize}

\subsection{BST Node Deletion: Three Cases}
\begin{enumerate}
    \item \textbf{Node is a Leaf:} Directly unlink and delete ($O(1)$).
    \item \textbf{Node has One Child:} Bypass node, linking parent directly to child ($O(1)$).
    \item \textbf{Node has Two Children:} Replace value with its \textbf{Inorder Successor} (smallest node in right subtree) or \textbf{Inorder Predecessor} (largest in left subtree), then recursively delete that successor/predecessor node ($O(h)$).
\end{enumerate}

\section{AVL Tree (Adelson-Velsky and Landis)}

An \textbf{AVL tree} is a strictly height-balanced BST where every node satisfies:
\[
\boxed{\text{Balance Factor (BF)} = h(\text{Left Subtree}) - h(\text{Right Subtree}) \in \{-1, \; 0, \; +1\}}
\]
If $|BF| \geq 2$, the tree is rebalanced using rotations:
\begin{table}[htbp]
\centering
\small
\begin{tabularx}{\textwidth}{l X X}
\toprule
\textbf{Imbalance Case} & \textbf{Triggering Insertion Location} & \textbf{Rebalancing Rotation} \\
\midrule
LL (Left-Left)   & Inserted in left subtree of left child   & Single Right Rotation \\
RR (Right-Right) & Inserted in right subtree of right child & Single Left Rotation \\
LR (Left-Right)  & Inserted in right subtree of left child  & Left Rotation on child, then Right on parent \\
RL (Right-Left)  & Inserted in left subtree of right child  & Right Rotation on child, then Left on parent \\
\bottomrule
\end{tabularx}
\caption{The four AVL tree imbalance scenarios and rotation solutions.}
\end{table}

\subsection{Minimum Nodes in an AVL Tree of Height $h$}
Let $N(h)$ be the minimum node count for an AVL tree of height $h$:
\[
\boxed{N(h) = 1 + N(h - 1) + N(h - 2)}, \quad N(-1) = 0, \; N(0) = 1
\]
This recurrence connects to the Fibonacci sequence:
\[
\boxed{N(h) = F_{h+3} - 1} \quad (F_0=0, F_1=1, F_2=1, F_3=2, \dots)
\]
\begin{center}
\small
\begin{tabular}{ccccccccc}
\toprule
$h$ & 0 & 1 & 2 & 3 & 4 & 5 & 6 \\
\midrule
$N(h)$ & 1 & 2 & 4 & 7 & 12 & 20 & 33 \\
\bottomrule
\end{tabular}
\end{center}

\section{Red-Black Tree}

A \textbf{Red-Black tree} is a self-balancing BST with node colors (Red/Black) ensuring:
\begin{enumerate}
    \item Every node is either Red or Black.
    \item The root is always Black.
    \item All external NIL leaves are Black.
    \item If a node is Red, both of its children must be Black (no two consecutive Reds).
    \item Every path from a node to descendant NIL leaves contains identical Black nodes (equal Black-height).
\end{enumerate}

\subsection{Height Bound}
For $n$ internal nodes, the height $h$ in edges satisfies:
\[
\boxed{h \leq 2 \log_2(n + 1)}
\]

\subsection{AVL vs Red-Black Tree Comparison}
\begin{itemize}
    \item \textbf{AVL:} More rigidly balanced. Faster lookups ($O(\log n)$ with smaller constant factors). Best for lookup-heavy workloads.
    \item \textbf{Red-Black:} More relaxed balance (height up to $2 \log_2 n$). Fewer rotations during insert/delete. Best for write-intensive workloads (used in C++ \texttt{std::map}, Java \texttt{TreeMap}, Linux CFS).
\end{itemize}

\begin{center}
\vspace{1em}
{\color{primary}\rule{0.45\textwidth}{1.2pt}}\\[4pt]
{\Large\bfseries\color{primary} Section IV: B-Trees and B+ Trees}\\[2pt]
{\color{primary}\rule{0.45\textwidth}{1.2pt}}
\vspace{0.5em}
\end{center}

\section{B-Tree}

A \textbf{B-tree} is a balanced, multiway search tree optimized for systems with block-based storage (disk pages, SSD blocks).

\subsection{Order $m$ Specification}
Let order $m$ denote the \textbf{maximum number of children} an internal node can have:
\begin{itemize}
    \item Maximum children per internal node: $m$.
    \item Maximum keys per node: $m - 1$.
    \item Minimum children in non-root internal node: $\lceil m/2 \rceil$.
    \item Minimum keys in non-root internal node: $\lceil m/2 \rceil - 1$.
    \item Root node: Has at least $2$ children (if non-leaf) and at least $1$ key.
    \item All leaf nodes reside at the exact same depth.
\end{itemize}

\section{What Is a B+ Tree?}

A \textbf{B+ tree} is a balanced multiway search tree where:
\begin{enumerate}
    \item \textbf{All actual data records or record pointers reside exclusively in leaf nodes.}
    \item Internal nodes store only \textbf{separator (routing) keys} and child pointers.
    \item All leaf nodes are linked sequentially in a \textbf{doubly linked list}, facilitating $O(1)$ contiguous range scans.
\end{enumerate}

\begin{table}[htbp]
\centering
\small
\begin{tabularx}{\textwidth}{l X X}
\toprule
\textbf{Feature} & \textbf{B-Tree} & \textbf{B+ Tree} \\
\midrule
Data Storage Location     & Internal nodes AND leaf nodes           & Leaf nodes exclusively \\
Internal Node Contents    & Keys, child pointers, AND data pointers & Separator routing keys and child pointers \\
Leaf Node Linking         & Leaves not linked together              & Leaves linked in sequential linked list \\
Range Queries             & Requires full tree inorder traversal    & Extremely fast sequential leaf scan \\
Redundant Keys            & Every key appears exactly once          & Keys in internal nodes duplicate leaf keys \\
Tree Height for same data & Taller (fewer keys fit per internal node)& Shorter (higher fanout per page) \\
Application               & File systems, search engines            & Database indexing (InnoDB, PostgreSQL, SQLite) \\
\bottomrule
\end{tabularx}
\caption{Structural comparison between B-Tree and B+ Tree architectures.}
\end{table}

\section{Critical B+ Tree Calculation: Maximum Number of Keys}

\subsection{Standard Convention Formulation}
Assuming:
\begin{itemize}
    \item Order $m$ = maximum children per internal node.
    \item Each internal node holds at most $m - 1$ separator keys.
    \item Each leaf holds at most $m - 1$ data records.
    \item Height $h$ = edge count from root to leaf (root at level $0$, leaves at level $h$).
\end{itemize}

The maximum number of leaf nodes at level $h$ is $m^h$. Therefore, the maximum number of data records is:
\[
\boxed{N_{\max} = (m - 1) m^h}
\]

\subsubsection{Example: Order $m=4$, Height $h=3$ Edges}
\begin{align*}
\text{Max leaf nodes} &= 4^3 = 64 \\
\text{Max data entries per leaf} &= 4 - 1 = 3 \\
\boxed{N_{\max}} &= 64 \times 3 = \mathbf{192 \text{ records}}
\end{align*}

\subsection{When Height Means Levels ($H$)}
If height is defined as total levels $H$ (where $h = H - 1$):
\[
\boxed{N_{\max} = (m - 1) m^{H - 1}}
\]
For $m=4, H=3$ levels: $N_{\max} = 3 \times 4^2 = 48$ records.

\section{Maximum Key Entries Across All B+ Tree Nodes}

If an exam question asks for \textbf{all physical key entries} stored across both internal and leaf nodes:
\begin{align*}
\text{Internal Nodes} &= \sum_{i=0}^{h-1} m^i = \frac{m^h - 1}{m - 1} \\
\text{Internal Separator Keys} &= (m - 1) \left( \frac{m^h - 1}{m - 1} \right) = m^h - 1 \\
\text{Leaf Data Keys} &= (m - 1) m^h \\
\text{Total Physical Key Slots} &= (m^h - 1) + (m - 1) m^h = \boxed{m^{h+1} - 1}
\end{align*}
For $m=4, h=3$: $K_{\text{total}} = 4^4 - 1 = 255$ key slots ($192$ data records $+$ $63$ internal separators).

\section{Minimum Number of Keys in a B+ Tree}

Let $a = \lceil m/2 \rceil$ be the minimum child capacity of non-root internal nodes, and $q = \lceil (m - 1)/2 \rceil$ be the minimum entry capacity of non-root leaves.

For height $h \geq 1$, the root has at least $2$ children, and subsequent internal nodes have at least $a$ children.
\[
\text{Min Leaf Nodes} = L_{\min} = 2 a^{h - 1} = 2 \left\lceil \frac{m}{2} \right\rceil^{h - 1}
\]
\[
\boxed{N_{\min} = L_{\min} \times q = 2 \left\lceil \frac{m}{2} \right\rceil^{h - 1} \left\lceil \frac{m - 1}{2} \right\rceil}
\]
For $m=4, h=3$:
\begin{itemize}
    \item $a = \lceil 4/2 \rceil = 2$, $q = \lceil 3/2 \rceil = 2$.
    \item $L_{\min} = 2 \times 2^{3-1} = 8$ leaves.
    \item $\boxed{N_{\min} = 8 \times 2 = 16 \text{ records}}$.
\end{itemize}

\section{B-Tree Key-Count Calculations}

In a B-tree of order $m$ and height $h$ edges, every node (internal and leaf) holds actual data keys:
\[
\boxed{K_{\max} = m^{h+1} - 1}
\]
\[
\boxed{K_{\min} = 2 \left\lceil \frac{m}{2} \right\rceil^h - 1}
\]
For $m=4, h=3$:
\[
K_{\max} = 4^4 - 1 = 255, \quad K_{\min} = 2(2^3) - 1 = 15 \text{ keys}
\]

\begin{center}
\vspace{1em}
{\color{primary}\rule{0.45\textwidth}{1.2pt}}\\[4pt]
{\Large\bfseries\color{primary} Section V: Other Important Tree Structures}\\[2pt]
{\color{primary}\rule{0.45\textwidth}{1.2pt}}
\vspace{0.5em}
\end{center}

\section{Trie (Prefix Tree)}
A tree structure used for fast string retrieval by sharing common prefixes.
\begin{itemize}
    \item \textbf{Height:} $h \leq L$, where $L$ is the length of the longest string.
    \item \textbf{Lookup Complexity:} $O(L)$ time, completely independent of the number of words $n$.
    \item \textbf{Max Nodes for alphabet $\sigma$ and length $L$:}
    \[
    N_{\max} = \frac{\sigma^{L+1} - 1}{\sigma - 1}
    \]
\end{itemize}

\section{Expression Trees}
Binary trees where leaves are operands and internal nodes are operators ($+,-,\times,\div$).
\begin{itemize}
    \item Preorder Traversal yields \textbf{Prefix notation}.
    \item Inorder Traversal yields \textbf{Infix notation}.
    \item Postorder Traversal yields \textbf{Postfix notation}.
    \item If all operators are binary: $L = I + 1 \implies \mathbf{n = 2I + 1}$.
\end{itemize}

\section{Huffman Coding Tree}
A full binary tree used for variable-length prefix-free data compression:
\begin{itemize}
    \item Internal Nodes: $I = L - 1$; Total Nodes: $n = 2L - 1$.
    \item Weighted Path Length (Cost): $C = \sum_{i=1}^L f_i d_i$, where $f_i$ is frequency and $d_i$ is tree depth.
\end{itemize}

\section{Segment Tree}
Stores interval information (range sums, minimums, maximums):
\begin{itemize}
    \item Full binary tree with $n$ leaves has exactly $\mathbf{2n - 1}$ logical nodes.
    \item Standard array allocation size is bounded by $\mathbf{4n}$.
    \item Build: $O(n)$; Range Query: $O(\log n)$; Point Update: $O(\log n)$.
\end{itemize}

\begin{center}
\vspace{1em}
{\color{primary}\rule{0.45\textwidth}{1.2pt}}\\[4pt]
{\Large\bfseries\color{primary} Section VI: Master Complexity Comparison}\\[2pt]
{\color{primary}\rule{0.45\textwidth}{1.2pt}}
\vspace{0.5em}
\end{center}

\begin{table}[htbp]
\centering
\small
\begin{tabularx}{\textwidth}{l X X X}
\toprule
\textbf{Tree Data Structure} & \textbf{Search Complexity} & \textbf{Insert Complexity} & \textbf{Delete Complexity} \\
\midrule
Ordinary Binary Tree & $O(n)$              & $O(n)$              & $O(n)$ \\
Ordinary BST         & $O(h)$ ($O(n)$ max)& $O(h)$              & $O(h)$ \\
Balanced BST (AVL)   & $O(\log n)$         & $O(\log n)$         & $O(\log n)$ \\
Red-Black Tree       & $O(\log n)$         & $O(\log n)$         & $O(\log n)$ \\
Binary Heap          & $O(n)$ arbitrary    & $O(\log n)$         & $O(\log n)$ extract root \\
B-Tree               & $O(\log n)$         & $O(\log n)$         & $O(\log n)$ \\
B+ Tree              & $O(\log n)$         & $O(\log n)$         & $O(\log n)$ \\
Trie (Prefix Tree)   & $O(L)$ key length   & $O(L)$              & $O(L)$ \\
Segment Tree         & $O(\log n)$ range   & Static build $O(n)$ & Point update $O(\log n)$ \\
Huffman Tree         & Static tree         & Build $O(n \log n)$ & Decode $O(L)$ \\
\bottomrule
\end{tabularx}
\caption{Master operational time complexity across tree architectures.}
\end{table}

\begin{center}
\vspace{1em}
{\color{primary}\rule{0.45\textwidth}{1.2pt}}\\[4pt]
{\Large\bfseries\color{primary} Section VII: Solved Mathematical Problems}\\[2pt]
{\color{primary}\rule{0.45\textwidth}{1.2pt}}
\vspace{0.5em}
\end{center}

\begin{quote}
\textbf{Problem 1: Nodes in Perfect Binary Tree of Height 6}\\
$n = 2^{6+1} - 1 = 2^7 - 1 = \mathbf{127 \text{ nodes}}$. Leaves $L = 2^6 = \mathbf{64}$. Internal $I = 64 - 1 = \mathbf{63}$.
\end{quote}

\begin{quote}
\textbf{Problem 2: Full Binary Tree with 25 Leaves}\\
$I = L - 1 = 25 - 1 = \mathbf{24 \text{ internal nodes}}$. Total $n = 24 + 25 = \mathbf{49 \text{ nodes}}$.
\end{quote}

\begin{quote}
\textbf{Problem 3: Minimum and Maximum Height for 40 Nodes}\\
$h_{\min} = \lceil \log_2(41) \rceil - 1 = 6 - 1 = \mathbf{5 \text{ edges}}$. $h_{\max} = n - 1 = \mathbf{39 \text{ edges}}$.
\end{quote}

\begin{quote}
\textbf{Problem 4: Leaves from Two-Child Node Count ($n=50, n_2=18$)}\\
$n_0 = n_2 + 1 = 18 + 1 = \mathbf{19 \text{ leaves}}$. Single child nodes $n_1 = 50 - 19 - 18 = \mathbf{13}$.
\end{quote}

\begin{quote}
\textbf{Problem 5: Can a Full Binary Tree Have Exactly 100 Nodes?}\\
No. A full binary tree has $n = 2I + 1$, which is strictly \textbf{odd}. Since $100$ is even, it is impossible.
\end{quote}

\begin{quote}
\textbf{Problem 6: Distinct BST Structures for 4 Keys}\\
$C_4 = \frac{1}{5} \binom{8}{4} = \frac{70}{5} = \mathbf{14 \text{ distinct BSTs}}$.
\end{quote}

\begin{quote}
\textbf{Problem 7: B+ Tree Capacity ($m=4$, Height 4 edges, 3 records/leaf)}\\
$L_{\max} = 4^4 = 256$ leaves. $\mathbf{N_{\max}} = 256 \times 3 = \mathbf{768 \text{ data records}}$.
\end{quote}

\begin{quote}
\textbf{Problem 8: B-Tree Keys ($m=6$, Height 2 edges)}\\
$K_{\max} = 6^{2+1} - 1 = 216 - 1 = \mathbf{215 \text{ keys}}$.\\
$K_{\min} = 2 \lceil 6/2 \rceil^2 - 1 = 2(3^2) - 1 = \mathbf{17 \text{ keys}}$.
\end{quote}

\begin{quote}
\textbf{Problem 9: Feasible Height for 18-Node AVL Tree}\\
$N(4) = 12, N(5) = 20$. Since $18 < 20$, height $5$ is impossible. Max feasible height is $\mathbf{h = 4}$.
\end{quote}

\begin{quote}
\textbf{Problem 10: Segment Tree Nodes for Array of Size 75}\\
$N = 2(75) - 1 = \mathbf{149 \text{ logical nodes}}$ (buffer bound $\approx 4 \times 75 = 300$).
\end{quote}

\begin{center}
\vspace{1em}
{\color{primary}\rule{0.45\textwidth}{1.2pt}}\\[4pt]
{\Large\bfseries\color{primary} Section VIII: Master Formula Cheat-Sheet}\\[2pt]
{\color{primary}\rule{0.45\textwidth}{1.2pt}}
\vspace{0.5em}
\end{center}

\begin{table}[htbp]
\centering
\small
\begin{tabularx}{\textwidth}{l X}
\toprule
\textbf{Tree Concept} & \textbf{Master Formula} \\
\midrule
Edge count for $n$ nodes & $E = n - 1$ \\
Max nodes at level $i$ & $m^i$ (general) $\;/\; 2^i$ (binary) \\
Max nodes for height $h$ & $\frac{m^{h+1}-1}{m-1}$ (general) $\;/\; 2^{h+1}-1$ (binary) \\
Min height for $n$ nodes & $\lceil \log_m(n(m-1)+1) \rceil - 1$ \\
Full binary tree leaves & $L = I + 1$ \\
Full binary tree total nodes & $n = 2I + 1$ (always odd) \\
Leaf vs 2-child relation & $n_0 = n_2 + 1$ (holds for \textbf{all} binary trees) \\
Null child pointers & $n + 1$ \\
Number of binary tree shapes & $C_n = \frac{1}{n+1} \binom{2n}{n}$ (Catalan Number) \\
AVL min nodes recurrence & $N(h) = 1 + N(h-1) + N(h-2) = F_{h+3}-1$ \\
Red-Black height upper bound & $h \leq 2 \log_2(n+1)$ \\
B+ Tree max data records & $(m - 1) m^h$ \\
B+ Tree max physical key slots & $m^{h+1} - 1$ \\
B+ Tree min data records ($h \geq 1$) & $2 \lceil m/2 \rceil^{h-1} \lceil (m-1)/2 \rceil$ \\
B-Tree max keys & $m^{h+1} - 1$ \\
B-Tree min keys ($h \geq 1$) & $2 \lceil m/2 \rceil^h - 1$ \\
Binary Heap height & $h = \lfloor \log_2 n \rfloor$ \\
Segment Tree logical nodes & $2n - 1$ \\
\bottomrule
\end{tabularx}
\caption{Comprehensive master mathematical formula sheet for recruitment examinations.}
\end{table}

\begin{center}
\vspace{1em}
{\color{primary}\rule{0.45\textwidth}{1.2pt}}\\[4pt]
{\Large\bfseries\color{primary} Section IX: Common Examination Traps}\\[2pt]
{\color{primary}\rule{0.45\textwidth}{1.2pt}}
\vspace{0.5em}
\end{center}

\begin{enumerate}
    \item \textbf{Edges vs Levels:} Height in edges ($h$) has $h + 1$ levels. If an exam defines height as levels $H$, replace $h$ with $H - 1$.
    \item \textbf{Full vs Complete:} Full means every node has $0$ or $2$ children. Complete means levels are fully packed with the last level filled left-to-right.
    \item \textbf{Even Node Count Rule:} A full binary tree \emph{cannot} have an even number of nodes.
    \item \textbf{BST Ordering vs Heap Ordering:} In a BST, left $<$ parent $<$ right (horizontal order). In a heap, parent $\geq$ children (vertical order, with no left/right constraint).
    \item \textbf{B+ Tree Separators vs Data Records:} Internal keys in a B+ tree are routing separators and duplicate leaf keys. Total physical key slots ($m^{h+1}-1$) exceeds unique data records ($(m-1)m^h$).
    \item \textbf{B+ Tree Order Ambiguity:} Always check whether order $m$ denotes maximum children ($m$ children, $m-1$ keys) or maximum keys ($2d$ keys, $2d+1$ children).
    \item \textbf{Arbitrary Key Search in Heap:} Searching for an arbitrary key in a binary heap takes $O(n)$, not $O(\log n)$.
\end{enumerate}

\begin{quote}
\textbf{Final Architectural Takeaway:} Tree structures balance operational time by exploiting geometric branching. Binary trees, AVL trees, and Red-Black trees optimize memory-resident logarithmic searches, while B-trees and B+ trees optimize disk and page block I/O by minimizing tree height through massive branching fanout.
\end{quote}
